\documentclass{standalone}
\usepackage{circuitikz}
\usetikzlibrary{calc}
\definecolor{circuitnotered}{HTML}{E5534B}
\definecolor{circuitnotemark}{HTML}{FE00FE}
\begin{document}
\begin{circuitikz}[american, line width=0.8pt]
\ctikzset{bipoles/length=1.2cm}
\foreach \x in {0,2,4,6,8,10} {\foreach \y in {0,-2,-4,-6,-8,-10,-12,-14,-16} {\fill[gray, opacity=0.35] (\x,\y) circle (1.1pt);}}
\node[gray, font=\scriptsize] at (0,0.84) {1};
\node[gray, font=\scriptsize] at (2,0.84) {2};
\node[gray, font=\scriptsize] at (4,0.84) {3};
\node[gray, font=\scriptsize] at (6,0.84) {4};
\node[gray, font=\scriptsize] at (8,0.84) {5};
\node[gray, font=\scriptsize] at (10,0.84) {6};
\node[gray, font=\scriptsize, anchor=east] at (-0.84,0) {a};
\node[gray, font=\scriptsize, anchor=east] at (-0.84,-2) {b};
\node[gray, font=\scriptsize, anchor=east] at (-0.84,-4) {c};
\node[gray, font=\scriptsize, anchor=east] at (-0.84,-6) {d};
\node[gray, font=\scriptsize, anchor=east] at (-0.84,-8) {e};
\node[gray, font=\scriptsize, anchor=east] at (-0.84,-10) {f};
\node[gray, font=\scriptsize, anchor=east] at (-0.84,-12) {g};
\node[gray, font=\scriptsize, anchor=east] at (-0.84,-14) {h};
\node[gray, font=\scriptsize, anchor=east] at (-0.84,-16) {i};
\coordinate (a1) at (0,0);
\coordinate (a5) at (8,0);
\draw (a1) to[R, l_=$R_{1}$, a^=$10\,\mathrm{k}\Omega$] (a5); % line 2
\draw[circuitnotered] (4,-4) circle (0.9); % line 4
\path (4,-4.395) rectangle (6.667,-3.605); % line 5
\node[anchor=west, inner sep=0, circuitnotemark, font=\footnotesize] at (4,-4) {X}; % line 5
\draw[circuitnotered] (4,-8) circle (0.9); % line 6
\path (2.349,-8.395) rectangle (5.651,-7.605); % line 7
\node[anchor=west, inner sep=0, circuitnotemark, font=\footnotesize] at (4,-8) {X}; % line 7
\draw[circuitnotered] (4,-12) circle (0.9); % line 8
\path (1.164,-12.395) rectangle (4,-11.605); % line 9
\node[anchor=west, inner sep=0, circuitnotemark, font=\footnotesize] at (4,-12) {X}; % line 9
\path (4,-16.395) rectangle (6.625,-15.605); % line 10
\node[anchor=west, inner sep=0, circuitnotemark, font=\footnotesize\bfseries] at (4,-16) {X}; % line 10
\path (12,-5.263) rectangle (18.852,0.395); % line 11
\node[anchor=west, inner sep=0, circuitnotemark, font=\footnotesize] at (12,0) {X}; % line 11
\node[anchor=west, inner sep=0, circuitnotemark, font=\footnotesize] at (12,-0.325) {X}; % line 11
\node[anchor=west, inner sep=0, circuitnotemark, font=\footnotesize] at (12,-0.649) {X}; % line 11
\node[anchor=west, inner sep=0, circuitnotemark, font=\footnotesize] at (12,-0.974) {X}; % line 11
\node[anchor=west, inner sep=0, circuitnotemark, font=\footnotesize] at (12,-1.298) {X}; % line 11
\node[anchor=west, inner sep=0, circuitnotemark, font=\footnotesize] at (12,-1.623) {X}; % line 11
\node[anchor=west, inner sep=0, circuitnotemark, font=\footnotesize] at (12,-1.947) {X}; % line 11
\node[anchor=west, inner sep=0, circuitnotemark, font=\footnotesize] at (12,-2.272) {X}; % line 11
\node[anchor=west, inner sep=0, circuitnotemark, font=\footnotesize] at (12,-2.596) {X}; % line 11
\node[anchor=west, inner sep=0, circuitnotemark, font=\footnotesize] at (12,-2.921) {X}; % line 11
\node[anchor=west, inner sep=0, circuitnotemark, font=\footnotesize] at (12,-3.246) {X}; % line 11
\node[anchor=west, inner sep=0, circuitnotemark, font=\footnotesize] at (12,-3.57) {X}; % line 11
\node[anchor=west, inner sep=0, circuitnotemark, font=\footnotesize] at (12,-3.895) {X}; % line 11
\node[anchor=west, inner sep=0, circuitnotemark, font=\footnotesize] at (12,-4.219) {X}; % line 11
\node[anchor=west, inner sep=0, circuitnotemark, font=\footnotesize] at (12,-4.544) {X}; % line 11
\node[anchor=west, inner sep=0, circuitnotemark, font=\footnotesize] at (12,-4.868) {X}; % line 11
\end{circuitikz}
\end{document}
